% 3-2-assumptions.tex
%  Budget: 3pp.  The six formal assumptions from the approved proposal.
%  Source map: interpretability named as a still-open challenge -> oconnor_review_2025 ;
%   held-out benchmark protocol / published comparison -> lago_forecasting_2021 .
%  CRITICAL (HANDOFF section 9 + amendment 2026-09-02): assumption (4) is
%   challenged on BOTH readings; (3) data quality is what the train/serve
%   definition mismatch actually implicates; (5) is challenged ONLY by a
%   residual this project did not isolate. Do NOT blame the German market
%   for the serving-path mismatch.

\section{فرضیات بنیادی پژوهش}
\label{sec:3-2}

هر پژوهش تجربی بر فرض‌هایی استوار است که اگر ناگفته بمانند، خواننده نمی‌داند نتایج
تا کجا معتبرند. شش فرض زیر همان‌هایی است که در پیشنهاده‌ی مصوب این پایان‌نامه
صورت‌بندی شده است. آن‌ها را در همین‌جا کامل می‌آوریم و برای هر یک تصریح می‌کنیم که
شواهد این پژوهش آن را تأیید کرده‌اند یا نه. این تصریح، عقب‌نشینی از پیشنهاد نیست؛
بازگشتِ مستند به آن است، و \cref{sec:5-2} حلقه‌ی همین بازگشت را می‌بندد.

\subsection*{فرض یکم: ایستایی}

فرض می‌شود ساختار آماریِ سری قیمت \lr{—} یا دست‌کم آن بخشی از آن که مدل‌ها بر آن
تکیه می‌کنند \lr{—} در طول دوره‌ی آموزش و آزمون به‌قدر کافی پایدار است که الگوی
آموخته‌شده بر گذشته، بر آینده نیز برقرار بماند.

این فرض در صورت قویِ خود برقرار نیست، و پژوهش نیز به آن نیازی ندارد. میانگین سالانه‌ی
قیمت در دوره‌ی مطالعه از \lr{42.89} یورو بر مگاوات‌ساعت در سال ۲۰۱۲ به \lr{28.98} در سال ۲۰۱۶
افت می‌کند و سپس به \lr{34.20} در سال ۲۰۱۷ بازمی‌گردد؛ سطح قیمت آشکارا ثابت نیست. آنچه
پژوهش در عمل فرض می‌کند صورتِ ضعیف‌ترِ آن است: واسنجی مجدد روزانه بر پنجره‌ی غلتانِ
۱۰۹۲روزه (\cref{sec:3-5}) مدل را پیوسته با سطح جاری بازار هم‌تراز نگه می‌دارد، به‌طوری‌که
تنها ایستاییِ محلی لازم می‌شود، نه ایستاییِ سراسری. این تمایز در فصل پنجم اهمیت پیدا
می‌کند، آنجا که مدل‌ها بدون هیچ واسنجی مجدد به کار گرفته می‌شوند.

\subsection*{فرض دوم: در دسترس بودن داده}

فرض می‌شود داده‌های لازم برای پیش‌بینی \lr{—} قیمت‌های تاریخی و پیش‌بینی‌های روزپیشِ
بار و تولید تجدیدپذیر \lr{—} پیش از مبدأ پیش‌بینی در دسترس‌اند.

فرض یادشده برقرار است، و در \cref{sec:3-4} به قیدی آزموده ترجمه شده است: هر ستون ماتریس ویژگی
باید تاریخ در دسترس بودن خود را اعلام کند، وگرنه ساخت ماتریس متوقف می‌شود. با
این‌همه فرض دوم یک هزینه‌ی ثبت‌شده دارد. چند متغیر که ادبیات آن‌ها را مؤثر می‌داند
\lr{—} به‌ویژه قیمت سوخت \lr{—} در این پژوهش در دسترس نبودند، زیرا سری‌های مرجعشان
تجاری و دارای مجوزند. \cref{sec:3-3} این محدودیت و دلیل آن را ثبت می‌کند؛ آنچه اینجا اهمیت
دارد این است که فرض دوم برای مجموعه‌ی ویژگیِ به‌کاررفته برقرار است، نه برای هر
مجموعه‌ی ویژگیِ مطلوب.

\subsection*{فرض سوم: کیفیت داده}

فرض می‌شود داده‌های ورودی از خطای اندازه‌گیری، جاافتادگی و ناسازگاریِ تعریفی مبرا
هستند.

برای مجموعه‌داده بنچمارک این فرض برقرار است: \lr{52{,}416} سطر ساعتی، بدون هیچ مقدار
گم‌شده و بدون هیچ برچسب زمانی تکراری. اما برای مجموعه‌ی داده‌ی زنده \lr{—} که تنها در
آزمون برون‌توزیعی به کار می‌رود \lr{—} این فرض نقض شده است، و نقض آن یکی از یافته‌های
همین پژوهش است. متغیر برون‌زای نخست در دوره‌ی آموزش پیش‌بینی بارِ ناحیه‌ی کنترلیِ
\lr{Amprion} است و در دوره‌ی سرویس‌دهی پیش‌بینی بارِ ملیِ آلمان: دو تعریف متفاوت زیر
یک نام. اندازه‌ی این ناهم‌خوانی سنجیده شده و در \cref{sec:5-2} گزارش می‌شود.

تصریح این نکته در همین‌جا اهمیت دارد، چون تعیین می‌کند شکست برون‌توزیعیِ فصل پنجم را
باید به کدام فرض نسبت داد. ناهم‌خوانیِ تعریفی، نقصی در مسیر داده‌ی خودِ ماست و ذیل
همین فرض سوم می‌نشیند، نه ذیل فرض پنجم.

\subsection*{فرض چهارم: تعمیم‌پذیری مدل}

فرض می‌شود مدلی که بر داده‌ی تاریخی آموزش دیده، بر داده‌ای که ندیده است نیز عملکرد
قابل‌قبولی دارد.

درون دوره‌ی بنچمارک این فرض برقرار است: دوره‌ی آزمونِ دوساله‌ی ۷۲۸ مبدأیی به‌طور کامل
کنار گذاشته شده و هیچ مدلی در هیچ مرحله‌ای \lr{—} نه در برازش، نه در تنظیم
ابرپارامتر و نه در برازش وزن‌های ترکیبی \lr{—} آن را ندیده است.

بیرون از آن دوره، این فرض برقرار نیست. مدل‌های تثبیت‌شده بر دوره‌ی بنچمارک، هنگام
ارزیابی بر ۱۷۳ روز داده‌ی زنده‌ی سال ۲۰۲۶، همگی از مدل ساده‌ی مرجع عقب می‌مانند و ترتیب
آن‌ها وارونه می‌شود. این صریح‌ترین نقضِ فرض در سرتاسر پژوهش است، و نکته‌ی مهم آنکه بر
هر دو خوانشِ ممکن برقرار می‌ماند: چه بخشی از افت را به تغییر بازار نسبت دهیم و چه به
ناهم‌خوانیِ تعریفیِ فرض سوم، در هر دو حالت مدل‌ها در شرایطی که برای آن آموزش ندیده
بودند شکست خورده‌اند. \cref{sec:5-2} اعداد آن را می‌آورد.

\subsection*{فرض پنجم: پایداری شرایط بازار}

فرض می‌شود شرایط کلانِ بازار میان دوره‌ی آموزش و دوره‌ی کاربرد به‌قدر کافی پایدار
می‌ماند. این فرض درباره‌ی \emph{بازار} است، و همین آن را از فرض چهارم جدا می‌کند.

نسبت‌دادنِ کاملِ شکست برون‌توزیعی به این فرض، خطایی است که این پژوهش عامدانه از آن
پرهیز می‌کند. بخشی از آن افت \lr{—} و بر پایه‌ی سنجش، بخش عمده‌ی سهمِ قابل اندازه‌گیریِ
آن \lr{—} از ناهم‌خوانیِ تعریفیِ مسیر سرویس‌دهیِ خودِ ما برمی‌خیزد، یعنی از فرض سوم،
و نه از حرکت بازار آلمان. آنچه با صداقت می‌توان گفت این است: فرض پنجم تنها با آن
باقی‌مانده‌ای به چالش کشیده می‌شود که پس از کنارگذاشتن سهم ناهم‌خوانی برجای می‌ماند،
و این پژوهش آن باقی‌مانده را جدا نکرده است. سهم بازار در این افت \lr{—} برخلاف سهم
ناهم‌خوانی \lr{—} به‌طور مستقل کمّی نشده است، و \cref{sec:5-2} آن را به همین صورت گزارش
می‌کند، نه صفر و نه غالب.

\subsection*{فرض ششم: تفسیرپذیری مدل}

فرض می‌شود می‌توان سهم ویژگی‌ها در خروجی مدل را به‌صورتی قابل‌فهم استخراج کرد.

در پیشنهاده‌ی مصوب، تفسیرپذیری یک فرض بود؛ در این پژوهش به یک تحویل‌دادنی بدل شده
است. مرورهای اخیر نیز تفسیرپذیری را در کنار کیفیت داده و تعمیم‌پذیری از چالش‌های
همچنان گشوده‌ی این حوزه برمی‌شمارند \cite{oconnor_review_2025}. \cref{sec:4-6} تحلیل
\lr{SHAP} را با دو قید صریح ارائه می‌کند: این تحلیل تنها بازوی درختی مدل را پوشش
می‌دهد و درباره‌ی اطلاعاتی که مدل خطی یا شبکه‌ی بازگشتی به کار می‌گیرند ادعایی نمی‌کند،
و مدلی که تفسیر می‌شود همزادِ وفادارِ مدل گزارش‌شده است و نه خودِ آن. فرض ششم با این
دو قید برقرار است.

\subsection*{جمع‌بندی}

از شش فرض بالا، فرض‌های دوم و ششم بی‌قید برقرارند؛ فرض یکم تنها در صورت ضعیف و محلیِ
خود لازم و برقرار است؛ فرض سوم برای داده‌ی بنچمارک برقرار و برای داده‌ی زنده نقض شده
است؛ فرض چهارم درون دوره‌ی بنچمارک برقرار و بیرون از آن نقض شده است؛ و فرض پنجم تنها
تا اندازه‌ای به چالش کشیده می‌شود که این پژوهش آن را از فرض سوم جدا نکرده است. \cref{sec:5-2} با شواهد عددی به همین فهرست بازمی‌گردد.
